\documentclass[conference]{IEEEtran}
\IEEEoverridecommandlockouts

\usepackage[T1]{fontenc}
\usepackage[utf8]{inputenc}
\usepackage[brazilian]{babel}
\usepackage{cite}
\usepackage{amsmath,amssymb,amsfonts}
\usepackage{graphicx}
\usepackage{textcomp}
\usepackage{xcolor}
\usepackage{url}

\renewcommand{\IEEEkeywordsname}{Palavras-chave}

\begin{document}

\title{Operational Models para Verifica\c{c}\~ao\\
de Brokers MQTT com ESBMC}

\author{
\IEEEauthorblockN{Gustavo Lima}
\IEEEauthorblockA{email@exemplo.org}
\and
\IEEEauthorblockN{Willian Jean}
\IEEEauthorblockA{email@exemplo.org}
\and
\IEEEauthorblockN{Paulo Gomes}
\IEEEauthorblockA{email@exemplo.org}
}

\maketitle

\begin{abstract}
Brokers MQTT s\~ao componentes cr\'iticos em arquiteturas de Internet das Coisas (IoT) e, por serem frequentemente escritos em linguagem~C, est\~ao sujeitos a vulnerabilidades de mem\'oria e aritm\'eticas. O \emph{Efficient SMT-Based Context-Bounded Model Checker} (ESBMC) \'e uma ferramenta de verifica\c{c}\~ao formal capaz de detectar tais defeitos, mas sua aplica\c{c}\~ao a software de rede \'e limitada pela aus\^encia de \emph{operational models} para APIs de sockets POSIX. Este trabalho prop\~oe o desenvolvimento desses modelos e sua integra\c{c}\~ao ao ESBMC, tendo o Eclipse Mosquitto como estudo de caso, com o objetivo de viabilizar a detec\c{c}\~ao automatizada de vulnerabilidades em brokers MQTT.
\end{abstract}

\begin{IEEEkeywords}
verifica\c{c}\~ao formal, model checking, ESBMC, MQTT, Mosquitto, IoT.
\end{IEEEkeywords}

\section{Introdu\c{c}\~ao}

\subsection{Contexto}
\label{sec:contexto}

O protocolo MQTT (\emph{Message Queuing Telemetry Transport}) \'e o padr\~ao de fato para comunica\c{c}\~ao em ambientes de Internet das Coisas (IoT), devido ao seu modelo leve de publica\c{c}\~ao/assinatura~\cite{mqtt_v5_spec}. O componente central de uma rede MQTT \'e o \emph{broker}, que recebe e roteia todas as mensagens entre dispositivos. O Eclipse Mosquitto~\cite{mosquitto} \'e uma das implementa\c{c}\~oes de broker mais adotadas, mantida pela Eclipse Foundation e escrita em linguagem~C.

Por ser escrito em C, o Mosquitto herda os riscos inerentes \`a linguagem: ger\^encia manual de mem\'oria, possibilidade de \emph{buffer overflows} e comportamentos indefinidos. Essas caracter\'isticas tornam software de rede escrito em~C particularmente suscet\'ivel a vulnerabilidades de seguran\c{c}a.

O ESBMC (\emph{Efficient SMT-Based Context-Bounded Model Checker})~\cite{cordeiro2012tse} \'e uma ferramenta de \emph{Bounded Model Checking} que verifica programas C/C++ traduzindo propriedades de seguran\c{c}a para f\'ormulas em l\'ogica SMT. O ESBMC \'e capaz de detectar automaticamente classes de defeitos como \emph{buffer overflow}, \emph{null pointer dereference}, estouro aritm\'etico e vazamentos de mem\'oria.

\subsection{Problema}
\label{sec:problema}

A aplica\c{c}\~ao do ESBMC a software de rede como o Mosquitto \'e limitada pela aus\^encia de \emph{operational models} para APIs externas usadas nesse dom\'inio. APIs de sockets POSIX (\texttt{socket}, \texttt{bind}, \texttt{recv}, \texttt{send}) e de multiplexa\c{c}\~ao de E/S (\texttt{select}, \texttt{poll}) n\~ao possuem modelos abstratos no ESBMC. Sem esses modelos, a ferramenta aborta a verifica\c{c}\~ao por s\'imbolos indefinidos ou tenta seguir a implementa\c{c}\~ao do sistema operacional, tornando a an\'alise invi\'avel.

O hist\'orico do Mosquitto registra CVEs que poderiam ser detectados por verifica\c{c}\~ao formal: \emph{integer overflow} no campo \emph{Remaining Length} (CVE-2017-7651), vazamento de mem\'oria no fluxo de conex\~ao (CVE-2017-7654) e \emph{off-by-one} em rotinas de parsing (CVE-2018-12551). Essas vulnerabilidades envolvem padr\~oes cl\'assicos de defeitos em~C que o ESBMC j\'a sabe detectar, mas a barreira de modelagem de APIs de rede impede a an\'alise.

O problema central \'e: como estender o ESBMC com \emph{operational models} para APIs de rede de modo a viabilizar a detec\c{c}\~ao automatizada de vulnerabilidades em brokers MQTT escritos em~C?

\subsection{Motiva\c{c}\~ao e Justificativa}
\label{sec:motivacao}

Brokers MQTT s\~ao pontos \'unicos de falha em arquiteturas IoT. O comprometimento de um broker permite interceptar ou injetar mensagens em toda a rede de dispositivos. Detectar vulnerabilidades antes da implanta\c{c}\~ao reduz o risco e o custo de opera\c{c}\~ao desses sistemas.

Do ponto de vista cient\'ifico, o cat\'alogo de \emph{operational models} do ESBMC para APIs de rede \'e incompleto. Preencher essa lacuna beneficia n\~ao apenas a an\'alise de brokers MQTT, mas toda uma classe de software de rede. Al\'em disso, os modelos s\~ao implementados em arquivos~C no diret\'orio de biblioteca do ESBMC, o que permite contribui\c{c}\~ao direta ao projeto \emph{open source} com baixo risco de regress\~ao.

A escolha do Mosquitto como estudo de caso se justifica por ser \emph{open source}, amplamente documentado e possuir um hist\'orico p\'ublico de CVEs que servem como or\'aculo de valida\c{c}\~ao.

\subsection{Objetivos}
\label{sec:objetivos}

\subsubsection{Objetivo Geral}

Desenvolver \emph{operational models} para APIs de rede no ESBMC e avaliar sua efic\'acia na detec\c{c}\~ao de vulnerabilidades no Eclipse Mosquitto.

\subsubsection{Objetivos Espec\'ificos}

\begin{enumerate}
    \item Identificar as APIs externas cuja aus\^encia de modelagem impede a verifica\c{c}\~ao do Mosquitto pelo ESBMC.

    \item Implementar \emph{operational models} para essas APIs no ESBMC.

    \item Construir \emph{harnesses} de verifica\c{c}\~ao para subsistemas cr\'iticos do Mosquitto, incluindo o parser MQTT e o tratamento do pacote \texttt{CONNECT}.

    \item Reproduzir CVEs hist\'oricos e recentes do Mosquitto com o ESBMC e comparar vers\~oes vulner\'aveis e corrigidas.

    \item Avaliar o impacto da extens\~ao em termos de tempo de verifica\c{c}\~ao, consumo de mem\'oria e taxa de detec\c{c}\~ao.
\end{enumerate}

\section{Metodologia}
\label{sec:metodologia}

A verifica\c{c}\~ao segue o padr\~ao de \emph{harnesses} isolados: trechos do c\'odigo Mosquitto (vers\~ao vulner\'avel ou corrigida) s\~ao copiados para arquivos C independentes, com estruturas reduzidas e entradas simb\'olicas (\texttt{\_\_VERIFIER\_nondet\_*}) ou pacotes concretos que disparam o defeito.

Para cada CVE, define-se uma propriedade de seguran\c{c}a expressa em \texttt{\_\_ESBMC\_assert} ou verificada por recursos nativos do ESBMC (\texttt{--unsigned-overflow-check}, \texttt{--memory-leak-check}, checagem de limites de \emph{array}). O ESBMC~8.3 executa \emph{bounded model checking} com limite de desenrolamento (\texttt{--unwind}).

As APIs de rede do broker completo n\~ao s\~ao verificadas de ponta a ponta nesta fase; o foco s\~ao subsistemas de parsing, ACL, propriedades MQTT~v5 e l\'ogica de ponteiros (bridge). Um cat\'alogo de APIs externas e \emph{stubs} m\'inimos est\~ao em \texttt{docs/api\_catalog.md} e \texttt{esbmc\_models/network\_stubs.c}.

\section{Cat\'alogo de CVEs do broker Mosquitto}
\label{sec:cves}

A Tabela~\ref{tab:cves} resume CVEs reproduzidos ou modelados. Vers\~oes atuais (2.1.x) incorporam as corre\c{c}\~oes de 2.0.16 e 2.0.19; n\~ao h\'a CVE p\'ublico novo listado para 2.1.x no ChangeLog oficial.

\begin{table}[htbp]
\caption{CVEs do broker Mosquitto usados como or\'aculo de valida\c{c}\~ao}
\label{tab:cves}
\centering
\begin{tabular}{llll}
\hline
CVE & Tipo & Corrigido em & Harness \\
\hline
2017-7651 & integer overflow & 1.4.15 & \texttt{remaining\_length} \\
2017-7654 & memory leak & 1.4.15 & \texttt{connect\_memleak} \\
2017-7650 & ACL bypass & 1.4.12 & \texttt{acl\_bypass} \\
2023-3592 & memory leak & 2.0.16 & \texttt{will\_properties} \\
2023-0809 & estado/protocolo & 2.0.16 & \texttt{initial\_packet} \\
2023-28366 & memory leak QoS2 & 2.0.16 & \texttt{qos2\_dup} \\
2024-3935 & double free & 2.0.19 & \texttt{bridge\_remap} \\
\hline
\end{tabular}
\end{table}

\section{Resultados}
\label{sec:resultados}

A Tabela~\ref{tab:results} resume execu\c{c}\~oes com ESBMC~8.3 (logs em \texttt{harnesses/mosquitto/*.log}). \textbf{FAILED} na vers\~ao vulner\'avel indica que a propriedade de seguran\c{c}a foi violada (detec\c{c}\~ao desejada). \textbf{SUCCESSFUL} na vers\~ao corrigida ou em propriedades do c\'odigo 2.1.x indica conformidade no limite de desenrolamento.

\begin{table}[htbp]
\caption{Resultados ESBMC nos harnesses (broker Mosquitto)}
\label{tab:results}
\centering
\scriptsize
\begin{tabular}{lll}
\hline
Harness & Flags principais & Resultado \\
\hline
\texttt{remaining\_length} (vuln.) & \texttt{--unsigned-overflow-check} & FAILED \\
\texttt{connect\_memleak} (vuln.) & \texttt{--memory-leak-check} & FAILED (CWE-401) \\
\texttt{read\_binary\_obo} (vuln.) & limites de \emph{array} & FAILED (CWE-193) \\
\texttt{acl\_bypass} (v1.4.10) & \texttt{--unwind 32} & FAILED \\
\texttt{acl\_bypass\_fixed} & \texttt{--unwind 32} & SUCCESSFUL \\
\texttt{will\_properties} (vuln.) & \texttt{--memory-leak-check} & FAILED \\
\texttt{will\_properties\_fixed} & \texttt{--memory-leak-check} & SUCCESSFUL \\
\texttt{initial\_packet} (vuln.) & --- & FAILED \\
\texttt{initial\_packet\_fixed} & --- & SUCCESSFUL \\
\texttt{qos2\_dup} (vuln.) & \texttt{--memory-leak-check} & FAILED \\
\texttt{qos2\_dup\_fixed} & \texttt{--unwind 8} & SUCCESSFUL \\
\texttt{bridge\_remap} (vuln.) & \texttt{--memory-leak-check} & FAILED \\
\texttt{bridge\_remap\_fixed} & \texttt{--unwind 16} & SUCCESSFUL \\
\texttt{varint\_harness} / \texttt{max\_packet\_size} & \texttt{--unwind 8} & SUCCESSFUL \\
\texttt{packet\_recv\_remaining\_length} (vuln.) & \texttt{network\_stubs.c} & FAILED \\
\texttt{packet\_recv\_remaining\_length\_fixed} & \texttt{network\_stubs.c} & SUCCESSFUL \\
\texttt{packet\_recv\_initial\_command} (vuln.) & \texttt{network\_stubs.c} & FAILED \\
\texttt{packet\_recv\_initial\_command\_fixed} & \texttt{network\_stubs.c} & SUCCESSFUL \\
\hline
\end{tabular}
\end{table}

Os quatro \emph{harnesses} de integra\c{c}\~ao (\texttt{packet\_recv\_*}) leem bytes via \texttt{recv()} usando o \emph{operational model} em \texttt{esbmc\_models/network\_stubs.c}, demonstrando detec\c{c}\~ao de CVE-2017-7651 e CVE-2023-0809 no fluxo de rede simb\'olica, n\~ao apenas com buffers fixos em mem\'oria.

O parser \texttt{packet\_\_read\_varint} do ramo 2.1.x satisfaz limites de \emph{Remaining Length} (m\'aximo quatro bytes e valor $\leq 2^{28}-1$), alinhado \`a especifica\c{c}\~ao MQTT~\cite{mqtt_v5_spec}.

\section{Limita\c{c}\~oes e trabalho futuro}
\label{sec:limitacoes}

A an\'alise \'e \emph{context-bounded}: o desenrolamento limita loops e pode omitir caminhos mais profundos. CVEs que envolvem sequ\^encias longas de pacotes (ex.: CVE-2024-8376) exigem modelagem de estado de sess\~ao e APIs \texttt{select}/\texttt{poll}, ainda n\~ao integradas ao broker completo.

Foram implementados \emph{operational models} m\'inimos (\texttt{network\_stubs.c}) e quatro \emph{harnesses} de integra\c{c}\~ao com \texttt{recv}. O pr\'oximo passo \'e incorporar modelos completos (\texttt{socket\_lib.c}) \`a biblioteca c2goto do ESBMC e verificar trechos reais de \texttt{packet\_mosq.c} do Mosquitto, medindo tempo e mem\'oria em rela\c{c}\~ao aos harnesses isolados.

\section{Conclus\~ao}
\label{sec:conclusao}

Este trabalho demonstra que o ESBMC detecta padr\~oes de vulnerabilidades reais do Eclipse Mosquitto---incluindo CVEs de 2017 e de 2023--2024---em harnesses isolados e em harnesses de integra\c{c}\~ao com \texttt{recv}, e que c\'odigo recente do parser MQTT satisfaz propriedades de seguran\c{c}a verificadas. A extens\~ao com modelos POSIX completos no ESBMC permanece como passo para verifica\c{c}\~ao do broker completo.

\bibliographystyle{IEEEtran}
\bibliography{refs}

\end{document}
